% Appendix A: proof of Proposition prop:Q (= JZC_M1 Theorem Q, archive/sec3.tex, proof unchanged).
\section{Proof of Proposition~\ref{prop:Q}}\label{app:Q}

By (H2), $g(y):=\sqrt y\,\sigma(y)=y\sigma(y)/\sqrt y$ is strictly decreasing and $\sigma(y)\ge y_0\sigma(y_0)/y$ for $y<y_0$; if $Y=\infty$, (H3)
gives $\int^\infty\sigma<\infty$ and $y\sigma(y)\to0$. Write $U^\chi(z)=\int\log|z-x|^{-1}\diff\chi(x)$.

\emph{The root.} Substituting $y=Eu$, $F(E)=\int_1^{Y/E}E\sigma(Eu)(1-1/u)^{-1/2}\diff u$. For fixed $u$, $E\sigma(Eu)=(Eu\,\sigma(Eu))/u$ is strictly
decreasing in $E$ by (H2), and the upper limit decreases, so $F$ is strictly decreasing; on $[E_0,E_1]\subset(0,Y)$, $E\sigma(Eu)\le E_0\sigma(E_0u)$, so
$F$ is continuous by dominated convergence. As $E\to0$, $F(E)\ge\int_E^Y\sigma\to\infty$. If $Y<\infty$, $F(E)\to0$ as $E\to Y$. If $Y=\infty$, split
at $2E$: $\int_{2E}^\infty\le\sqrt2\int_{2E}^\infty\sigma\to0$, and $\int_E^{2E}\le E\sigma(E)\int_1^2(1-1/u)^{-1/2}\diff u\to0$.

\emph{The density.} Let $S(z):=\int_E^Y\tilde\sigma(y)(z-y)^{-1}\diff y$, $\tilde G(z):=((z-E)/z)^{1/2}S(z)$ (principal branch) and $G(z):=\int(z-x)^{-1}\diff\chi(x)$.
For fixed $y\in(E,Y)$ put
\[
  k_y(z):=\frac{((z-E)/z)^{1/2}-\sqrt{(y-E)/y}}{z-y},\qquad \rho_y(x):=\frac1\pi\sqrt{\frac{E-x}x}\,\frac1{y-x}>0\quad(0<x<E).
\]
The singularity of $k_y$ at $y$ is removable, so $k_y$ is analytic off $[0,E]$, $O(|z|^{-1/2})$ at $0$, bounded at $E$, and
$\sim(1-\sqrt{(y-E)/y})/z$ at infinity; it jumps by $-2\pi i\rho_y$ on $(0,E)$. By Liouville's theorem $k_y(z)=\int_0^E\rho_y(x)(z-x)^{-1}\diff x$,
and $\rho_y$ has mass $1-\sqrt{(y-E)/y}$. Since $\sigma=\tilde\sigma\sqrt{(y-E)/y}$, Fubini's theorem (all terms positive, $|z-x|^{-1}\le|\Im z|^{-1}$,
$\int \tilde\sigma(1-\sqrt{(y-E)/y})\le F(E)$) gives
\[
  \tilde G(z)=\int_E^Y\tilde\sigma(y)\Bigl[\frac{\sqrt{(y-E)/y}}{z-y}+k_y(z)\Bigr]\diff y=\int_E^Y\frac{\sigma(y)\diff y}{z-y}+\int_0^E\frac{\psi(x)\diff x}{z-x}=G(z),
\]
because $\int \tilde\sigma(y)\rho_y(x)\diff y=\psi(x)$. Hence $\chi$ has mass $\int \tilde\sigma=F(E)=1$, $\Re G(x\pm i0)=0$ on $(0,E)$, and $\psi>0$. With $R(y):=\sqrt{y(y-E)}$,
$R\sigma=\tilde\sigma(y-E)$ gives $\tilde G=R(z)^{-1}[F(E)+\int\sigma R(z-y)^{-1}\diff y]$, so also
$\pi\sqrt{x(E-x)}\,\psi(x)=1-\int_E^Y\sigma(y)\sqrt{y(y-E)}(y-x)^{-1}\diff y$ by~\eqref{eq:Q}. For the upper bound write $\tilde\sigma=g\cdot(y-E)^{-1/2}$ and use
$\int_E^\infty(y-E)^{-1/2}(y-x)^{-1}\diff y=\pi(E-x)^{-1/2}$: $\psi(x)<g(E)/\sqrt x\le g(x)/\sqrt x=\sigma(x)$.

\emph{The potential.} On $(0,E)$, $U^\chi$ is constant because $\Re G=0$. For $z\in(E,Y)$, $\varphi:=U^\chi-\ell$ satisfies $\varphi(E)=0$ and
$\varphi'(z)=-((z-E)/z)^{1/2}H(z)$, $H(z):=\PV\!\int_E^Y\tilde\sigma(y)(z-y)^{-1}\diff y$, which exists since $\sigma$ is locally H\"older. With $\tilde\sigma=g(y-E)^{-1/2}$,
\[
  H(z)=\int_E^Y\frac{(g(y)-g(z))(y-E)^{-1/2}}{z-y}\diff y+g(z)\,\PV\!\!\int_E^Y\frac{(y-E)^{-1/2}}{z-y}\diff y ;
\]
the first integrand is positive since $g$ decreases, and the principal value is $0$ if $Y=\infty$ and
$(z-E)^{-1/2}\log\frac{\sqrt{Y-E}+\sqrt{z-E}}{\sqrt{Y-E}-\sqrt{z-E}}>0$ if $Y<\infty$. Hence $\varphi<0$ on $(E,Y]$.

\emph{Uniqueness.} Since $\log^+|x-y|\le\log(1+x)+\log(1+y)$ and (H3) bounds the logarithmic moment of every $\chi'\le\sigma$, $I(\chi')>-\infty$ for every
competitor. Let $\chi'\le\sigma$ have unit mass and $I(\chi')<\infty$, $\upsilon:=\chi'-\chi$. Then $I(\chi')-I(\chi)=I(\upsilon)+2\int(U^\chi-\ell)\diff\upsilon$. The
integrand vanishes on $[0,E]$ ($U^\chi$ is continuous, the density of $\chi$ being at most $g(E)x^{-1/2}$ on $(0,E)$ and bounded on $[E,Y]$); on
$(E,Y]$, $\upsilon\le0$ and $U^\chi-\ell<0$, so the integral is $\ge0$. The signed measure $\upsilon$ has mass $0$ and
$\iint|\log|x-y||\diff|\upsilon|\diff|\upsilon|<\infty$ (the positive part by (H3), the negative part by the density bound and $I(\chi')<\infty$), so
$I(\upsilon)\ge0$ with equality only if $\upsilon=0$, by \cite[Lemma~3.3]{PaperI}.\qed

\begin{remark}\label{rem:Qedge}
Under~\eqref{eq:Q}, $\psi(E^-)=\sigma(E)$, so $\chi$ is continuous at $E$; without~\eqref{eq:Q}, $\psi$ acquires a multiple of $(E-x)^{-1/2}$. If $\sigma$
is H\"older of order $>\frac12$ near $E$, then $\sigma-\psi=C\sqrt{E-x}(1+o(1))$ with $C=-h(E)/(\pi\sqrt E)>0$,
$h(E):=\int_E^Y(g(y)-g(E))(y-E)^{-3/2}\diff y-2g(E)(Y-E)^{-1/2}\mathbf 1_{Y<\infty}$. The density and potential steps use only that $g$ is
decreasing; (H2) enters through the monotonicity of $F$. The hypotheses are not necessary but cannot simply be dropped: in a discretised energy
minimisation, a constraint with a dip below the edge gives three roots of~\eqref{eq:Q} and two bands. The root, the two forms of $\psi$, the
mass and the inequalities were checked by quadrature for the lattice at $p=0.6$, $1.5$ and two further densities (script
\texttt{m1\_theoremQ}).
\end{remark}
